\centerline{\bf \huge 9 класс}
\centerline{\normalsize{\it Работа рассчитана на \bf{180} минут}}

\problem{9.1}
Найдите наибольшее пятизначное число, произведение цифр которого равно 120.

\problem{9.2}
Лентяйка Эвена жёстко готовится к муниципу по математике. Чтобы не растратить все силы, она решает не более 10 задач в день, а если в какой-нибудь день она решала больше 7 задач, то следующие два дня она решала не более 5 задач в день. Какое наибольшее количество задач Эвена могла решить за 7 подряд идущих дней?

\problem{9.3}
Про положительные числа $a, b, c$ известно, что

$$
\frac{a+b+c}{a+b-c}=7, \quad \frac{a+b+c}{a+c-b}=1,75 .
$$


Чему равняется $\frac{a+b+c}{b+c-a}$ ?

%все задачи выше муницип Москва 2020

\problem{9.4}
Алина и Маша разложили на столе 13 различных карт. Каждая карта может лежать в одном из двух положений: рубашкой вверх или рубашкой вниз. Игроки должны по очереди переворачивать по одной карте. Проигрывает тот игрок, после хода которого повторится какая-то из предыдущих ситуаций (включая изначальную). Первый ход сделала Маша. Кто сможет выиграть независимо от того, как будет играть соперник?

%Москва 2016

\problem{9.5}
В остроугольном треугольнике $A B C$ проведена медиана $B M$. Точки $P$ и $Q$ --- центры вписанных окружностей треугольников $A B M$ и $C B M$ соответственно. Докажите, что вторая точка пересечения описанных окружностей треугольников $A B P$ и $C B Q$ лежит на отрезке $B M$.

%Москва 2018